%%%%%% Run at command line, run
%%%%%% xelatex grad-sample.tex 
%%%%%% for a few times to generate the output pdf file
\documentclass[12pt,oneside,openright,a4paper]{cpe-thai-project}


\usepackage{polyglossia}
\setdefaultlanguage{thai}
\setotherlanguage{english}
\newfontfamily\thaifont[Script=Thai,Scale=1.23]{TH Sarabun New}
\defaultfontfeatures{Mapping=tex-text,Scale=1.23,LetterSpace=0.0}
\setmainfont[Scale=1.23,LetterSpace=0,WordSpace=1.0,FakeStretch=1.0,Mapping=tex-text]{TH Sarabun New}
\XeTeXlinebreaklocale "th"	
\XeTeXlinebreakskip = 0pt plus 0pt
\emergencystretch=10pt

%%%%%%%%%%%%%%%%%%%%%%%%%%%%%%%%%%%%%%%%%%%%%%%%%%%%%%%%%%%%%%%%%%%
% Customize below to suit your needs 
% The ones that are optional can be left blank. 
%%%%%%%%%%%%%%%%%%%%%%%%%%%%%%%%%%%%%%%%%%%%%%%%%%%%%%%%%%%%%%%%%%%
% First line of title
\def\disstitleone{Document management system for lawyer}   
% Second line of title
% \def\disstitletwo{Project/Indep title line 2 (optional)}   
% Your first name and lastname
\def\dissauthor{Mr. Kittipat Dechkul}   % 1st member
\def\dissauthortwo{Mr. Choolerk Taebanpakul}   % 2nd member (optional)
\def\dissauthorthree{}   % 2nd member (optional)


% The degree that you're persuing..
\def\dissdegree{Bachelor of Engineering} % Name of the degree
\def\dissdegreeabrev{B.Eng} % Abbreviation of the degree
\def\dissyear{2022}                   % Year of submission
\def\thaidissyear{2565}               % Year of submission (B.E.)

%%%%%%%%%%%%%%%%%%%%%%%%%%%%%%%%%%%%%%%%%%%%
% Your project and independent study committee..
%%%%%%%%%%%%%%%%%%%%%%%%%%%%%%%%%%%%%%%%%%%%
\def\dissadvisor{Assoc.Prof. Priyakorn Pusawiro , Ph.D.}  % Advisor
\def\disscoadvisor{}  % Co-advisor
\def\disscommitteetwo{Assoc.Prof. Sanan Srakaew, Ph.D.}  % 3rd committee member (optional)
\def\disscommitteethree{Asst.Prof. Suthathip Maneewongpathana, Ph.D.}   % 4th committee member (optional) 
\def\disscommitteefour{Mr. Kharittha Jangsamsee, Ph.D.}    % 5th committee member (optional) 

\def\worktype{Project} %%  Project or Independent study
\def\disscredit{3}   %% 3 credits or 6 credits


\def\fieldofstudy{Computer Engineering} 
\def\department{Computer Engineering} 
\def\faculty{Engineering}

\def\thaifieldofstudy{วิศวกรรมคอมพิวเตอร์} 
\def\thaidepartment{วิศวกรรมคอมพิวเตอร์} 
\def\thaifaculty{วิศวกรรมศาสตร์}
 
\def\appendixnames{Appendix} %%% Appendices or Appendix

\def\thaiworktype{ปริญญานิพนธ์} %  Project or research project % 
\def\thaidisstitleone{Document management system for lawyer}
% \def\thaidisstitletwo{หัวข้อปริญญานิพนธ์บรรทัดสอง}
\def\thaidissauthor{นายกิตติภัฎ เดชกุล}
\def\thaidissauthortwo{นายชูฤกษ์ แต่บรรพกุล} %Optional
\def\thaidissauthorthree{} %Optional

\def\thaidissadvisor{รศ.ดร. ปริยกร ปุสวิโร}
%% Leave this empty if you have no co-advisor
% \def\thaidisscoadvisor{นายสุรสิทธิ์ ประคุณหังสิต} %Optional
\def\thaidissdegree{วิศวกรรมศาสตรบัณฑิต}

% Change the line spacing here...
\linespread{1.15}

%%%%%%%%%%%%%%%%%%%%%%%%%%%%%%%%%%%%%%%%%%%%%%%%%%%%%%%%%%%%%%%%
% End of personal customization.  Do not modify from this part 
% to \begin{document} unless you know what you are doing...
%%%%%%%%%%%%%%%%%%%%%%%%%%%%%%%%%%%%%%%%%%%%%%%%%%%%%%%%%%%%%%%%


%%%%%%%%%%%% Dissertation style %%%%%%%%%%%
%\linespread{1.6} % Double-spaced  
%%\oddsidemargin    0.5in
%%\evensidemargin   0.5in
%%%%%%%%%%%%%%%%%%%%%%%%%%%%%%%%%%%%%%%%%%%
%\renewcommand{\subfigtopskip}{10pt}
%\renewcommand{\subfigbottomskip}{-5pt} 
%\renewcommand{\subfigcapskip}{-6pt} %vertical space between caption
%                                    %and figure.
%\renewcommand{\subfigcapmargin}{0pt}

\renewcommand{\topfraction}{0.85}
\renewcommand{\textfraction}{0.1}

\newtheorem{theorem}{Theorem}
\newtheorem{lemma}{Lemma}
\newtheorem{corollary}{Corollary}

\def\QED{\mbox{\rule[0pt]{1.5ex}{1.5ex}}}
\def\proof{\noindent\hspace{2em}{\itshape Proof: }}
\def\endproof{\hspace*{\fill}~\QED\par\endtrivlist\unskip}
%\newenvironment{proof}{{\sc Proof:}}{~\hfill \blacksquare}
%% The hyperref package redefines the \appendix. This one 
%% is from the dissertation.cls
%\def\appendix#1{\iffirstappendix \appendixcover \firstappendixfalse \fi \chapter{#1}}
%\renewcommand{\arraystretch}{0.8}
%%%%%%%%%%%%%%%%%%%%%%%%%%%%%%%%%%%%%%%%%%%%%%%%%%%%%%%%%%%%%%%%
%%%%%%%%%%%%%%%%%%%%%%%%%%%%%%%%%%%%%%%%%%%%%%%%%%%%%%%%%%%%%%%%

\usepackage{ragged2e}
\begin{document}

\pdfstringdefDisableCommands{%
\let\MakeUppercase\relax
}

\begin{center}
  \includegraphics[width=2.8cm]{logo02.jpg}
\end{center}
\vspace*{-1cm}

\maketitlepage
\makesignaturepage 

%%%%%%%%%%%%%%%%%%%%%%%%%%%%%%%%%%%%%%%%%%%%%%%%%%%%%%%%%%%%%%
%%%%%%%%%%%%%%%%%%%%%% English abstract %%%%%%%%%%%%%%%%%%%%%%%
%%%%%%%%%%%%%%%%%%%%%%%%%%%%%%%%%%%%%%%%%%%%%%%%%%%%%%%%%%%%%%
\abstract

In a multihop ad hoc network, the interference among nodes is
  reduced to maximize the throughput by using a smallest transmission
  range that still preserve the network connectivity. However, most
  existing works on transmission range control focus on the
  connectivity but lack of results on the throughput performance. This
  paper analyzes the per-node saturated throughput of an IEEE 802.11b
  multihop ad hoc network with a uniform transmission range. Compared
  to simulation, our model can accurately predict the per-node
  throughput.  The results show that the maximum achievable per-node
  throughput can be as low as 11\% of the channel capacity in a normal
  set of $\alpha$ operating parameters independent of node density. However, if
  the network connectivity is considered, the obtainable throughput
  will reduce by as many as 43\% of the maximum throughput. 

\begin{flushleft}
\begin{tabular*}{\textwidth}{@{}lp{0.8\textwidth}}
\textbf{Keywords}: & Web application / (DMS) Document Management System / (OCR) Optical Character Recognition
\end{tabular*}
\end{flushleft}
\endabstract

%%%%%%%%%%%%%%%%%%%%%%%%%%%%%%%%%%%%%%%%%%%%%%%%%%%%%%%%%%%%%%
%%%%%%%%%% Thai abstract here %%%%%%%%%%%%%%%%%%%%%%%%%%%%%%%%%
%%%%%%%%%%%%%%%%%%%%%%%%%%%%%%%%%%%%%%%%%%%%%%%%%%%%%%%%%%%%%%
% {\newfontfamily\thaifont{TH Sarabun New:script=thai}[Scale=1.3]
% \XeTeXlinebreaklocale "th_TH"	
% \thaifont
\thaiabstract

เซ็นเซอร์ เอ็กซ์เพรสรองรับคอนเซปต์สหัสวรรษเมจิก อิ่มแปร้ เฟรชชี่ ชาร์ปเช็งเม้งคลาสสิก แพตเทิร์น แอลมอนด์ เพลซว้อยก๊วน ซาร์ดีนซี้เนิร์สเซอรีอีสต์ สเตเดียมเพียบแปร้โอ้ยแคมปัส จัมโบ้ช็อตแมคเคอเรลอึ๋ม สตริง แมกกาซีนสตริงผ้าห่ม ฮัลโหล ยิม รอยัลตี้

\begin{flushleft}
\begin{tabular*}{\textwidth}{@{}lp{0.8\textwidth}}
 & \\

\textbf{คำสำคัญ}: & การชุบเคลือบด้วยไฟฟ้า / การชุบเคลือบผิวเหล็ก /  เคลือบผิวรังสี
\end{tabular*}
\end{flushleft}
\endabstract

%}

%%%%%%%%%%%%%%%%%%%%%%%%%%%%%%%%%%%%%%%%%%%%%%%%%%%%%%%%%%%%
%%%%%%%%%%%%%%%%%%%%%%% Acknowledgments %%%%%%%%%%%%%%%%%%%%
%%%%%%%%%%%%%%%%%%%%%%%%%%%%%%%%%%%%%%%%%%%%%%%%%%%%%%%%%%%%
\preface
ขอบคุณอาจารย์ที่ปรึกษา กรรมการ พ่อแม่พี่น้อง และเพื่อนๆ คนที่ช่วยให้งานสำเร็จ ตามต้องการ

%%%%%%%%%%%%%%%%%%%%%%%%%%%%%%%%%%%%%%%%%%%%%%%%%%%%%%%%%%%%%
%%%%%%%%%%%%%%%% ToC, List of figures/tables %%%%%%%%%%%%%%%%
%%%%%%%%%%%%%%%%%%%%%%%%%%%%%%%%%%%%%%%%%%%%%%%%%%%%%%%%%%%%%
% The three commands below automatically generate the table 
% of content, list of tables and list of figures
\tableofcontents                    
\listoftables
\listoffigures                      

%%%%%%%%%%%%%%%%%%%%%%%%%%%%%%%%%%%%%%%%%%%%%%%%%%%%%%%%%%%%%%
%%%%%%%%%%%%%%%%%%%%% List of symbols page %%%%%%%%%%%%%%%%%%%
%%%%%%%%%%%%%%%%%%%%%%%%%%%%%%%%%%%%%%%%%%%%%%%%%%%%%%%%%%%%%%
% You have to add this manually..
\listofsymbols
\begin{flushleft}
\begin{tabular}{@{}p{0.07\textwidth}p{0.7\textwidth}p{0.1\textwidth}}
\textbf{SYMBOL}  & & \textbf{UNIT} \\[0.2cm]
% $\alpha$ & Test variable\hfill & m$^2$ \\
% $\lambda$ & Interarival rate\hfill &  jobs/second\\
% $\mu$ & Service rate\hfill & jobs/second\\
\end{tabular}
\end{flushleft}
%%%%%%%%%%%%%%%%%%%%%%%%%%%%%%%%%%%%%%%%%%%%%%%%%%%%%%%%%%%%%%
%%%%%%%%%%%%%%%%%%%%% List of vocabs & terms %%%%%%%%%%%%%%%%%
%%%%%%%%%%%%%%%%%%%%%%%%%%%%%%%%%%%%%%%%%%%%%%%%%%%%%%%%%%%%%%
% You also have to add this manually..
\listofvocab
\begin{flushleft}
\begin{tabular}{@{}p{1in}@{=\extracolsep{0.5in}}l}
  OCR &  Optical Character Recognition \\
\end{tabular}
\end{flushleft}

%\setlength{\parskip}{1.2mm}

%%%%%%%%%%%%%%%%%%%%%%%%%%%%%%%%%%%%%%%%%%%%%%%%%%%%%%%%%%%%%%%
%%%%%%%%%%%%%%%%%%%%%%%% Main body %%%%%%%%%%%%%%%%%%%%%%%%%%%%
%%%%%%%%%%%%%%%%%%%%%%%%%%%%%%%%%%%%%%%%%%%%%%%%%%%%%%%%%%%%%%%


\chapter{บทนำ}

\section{ที่มาและความสำคัญ}
\hspace*{1cm}ในปัจจุบัน ได้เกิด คดีต่างๆมากมาย จำนวนของ ทนายนั้นมีไม่เพียงพอ ในการ จัดการกับแต่ละคดี เลยจำเป็นที่ทนาย 1 คนจะต้องดูแล ในหลายคดีพร้อมๆ กัน ซึ่งในแต่ละคดีนั้น ทนายต้องเป็นคน จัดการกับเอกสารมากมาย ซึ่ง ทนายต้องจำให้ได้ว่า เอกสารไหน ที่ใช้กับคดีของคนไหน ซึ่งอาจจะ เกิดข้อผิด พลาด หรือเกิดความล่าช้า การการดำเนินการ ในแต่ละคดีได้ ดังนั้น จึงอยากทำระบบ DMS เพื่อให้ทนายนั้น สามารถ ทำงานได้มีประสิทธิภาพมากยิ่งขึ้น โดยมีประชาชนทั่วไปที่ต้องการจะติดตามรูปคดีของคดีต่างๆที่เกิดขึ้นนั้น สามารถเข้าถึงได้ยาก หรือไม่ก็ต้องการที่จะดูเอกสารของเคสต่างๆที่คนทั่วไปสามารถเข้าดูได้ แต่ว่าข้อมูลเหล่านี้เข้าถึงได้ยาก หรือบางคนก็ไม่รู้คำว่าจะต้องค้นหาอย่างไร \newline
\hspace*{1cm}ด้วยเหตุเหล่านี้ทางคณะผู้จัดทำได้เล็งเห็นถึงปัญหาดังกล่าว จึงต้องการที่จะพัฒนาเว็บแอปพลิเคชันเพื่อใช้ในการจัดการเอกสารต่างๆของทนาย โดยสามารถที่จะนำเอกสารบางส่วนที่สามารถเผยแพร่ได้ ให้ทุกคนสามารถเข้าถึงได้ และนำข้อมูลต่างๆมาแสดงผลในรูปแบบของแผนที่หรือแผนภูมิรูปต่างๆ

\section{วัตถุประสงค์}
\begin{enumerate}
  \item เพื่อศึกษาระบบในการบริหารจัดการเอกสาร การติดตาม และขั้นตอนการดำเนินคดีต่างๆ ของสำนักทนายความ
  \item เพื่อออกแบบฐานข้อมูล ให้เหมาะสมหรือสอดคล้องกับการค้นหาและจัดเก็บเอกสารอิเล็กทอนิค
  \item เพื่อพัฒนาเว็บแอปพลิเคชันที่ใช้ในการบริหารจัดการเอกสาร
  \item เพื่อสร้างช่องทางเข้าถึงข้อมูลเบื้องต้นอย่างง่ายให้กับทนายความและลูกความ
\end{enumerate}

\section{ขอบเขตและข้อจำกัดของปริญญานิพนธ์}
\hspace*{1cm}เว็บแอปพลิเคชันของคณะผู้จัดทำสามารถทำงานได้บนคอมพิวเตอร์ผ่านเว็บเบราว์เซอร์ต่าง ๆ เท่านั้นอาทิเช่น Google Chrome, Safari หรือ Microsoft Edge อีกทั้งข้อมูลที่นำมาใช้ในปริญญานิพนธ์นี้จะเป็นข้อมูลจริงที่ถูกเผยแพร่ผ่านเว็บไซต์ของศาลยุติธรรมซึ่งข้อมูลนี้ทุกคนสามารถเข้าถึงข้อมูลเหล่านี้ได้ และอีกส่วนหนึ่งเป็นข้อมูลเอกสารที่ได้ความยินยอมจากเจ้าของคดีและทนายให้สามารถเผยแพร่และเป็นตัวอย่างในปริญญานิพนธ์นี้เท่านั้น

\section{เนื้อหาทางวิศวกรรมที่เป็นต้นฉบับ}
\hspace*{1cm}โครงงานนี้พัฒนาขึ้นจากการใช้ความรู้ในด้านการออกแบบฐานข้อมูล ความรู้ในการพัฒนาเว็บแอปพลิเคชัน รวมไปถึงความรู้ทางด้านวิศวกรรมซอฟต์แวร์ ซึ่งไม่ว่าจะเป็นการออกแบบ UML diagram ต่างๆ การวางแผนการทำงาน การพัฒนาซอฟต์แวร์ในรูปแบบของ Waterfall การใช้ Kanbann เพื่อช่วยในการจัดการการทำงาน มาใช้เพื่อสร้างเป็นเว็ปแอปพลิเคชันสำหรับจัดการเอกสารของทนาย โดยมีการใช้ Elasticsearch ในการค้นหาเอกสาร และใช้ Gin framework ที่ใช้ภาษา Go language มาจัดการส่วนของการประมวณผล สุดท้ายได้นำ NextJs Framework ซึ่งใช้ภาษา Typescript , component framework และ CSS framework มาใช้ในการ แสดงข้อมูลที่ถูกประมวลผลแล้ว 

\newpage
\section{ขั้นตอนการทํางานและระยะเวลาการดําเนินงาน}
\begin{enumerate}
  \item เลือกหัวข้อที่สนใจ
  \item ศึกษาหัวข้อปริญญานิพนธ์ที่ได้รับ
  \item จัดทำข้อเสนอหัวข้อปริญญานิพนธ์
  \item ศึกษาทฤษฎี เทคโนโลยีและซอฟท์แวร์ที่ใช้
  \item เก็บโครงสร้างของข้อมูลที่ได้รับจากอุปกรณ์เครือข่าย
  \item วางโครงสร้างและรูปแบบของแอปพลิเคชั่น
  \item จัดทำรูปเล่มปริญญานิพนธ์ประจำภาคการศึกษาที่ 1/2565
  \item นำเสนอปริญญานิพนธ์ประจำภาคการศึกษาที่ 1/2565
  \item พัฒนาเว็บแอปพลิเคชัน
  \item ทดสอบเว็บแอปพลิเคชัน
  \item ปรับปรุงและแก้ไขเว็บแอปพลิเคชัน
  \item ทดสอบการใช้งานจริง
  \item จัดทำรูปเล่มปริญญานิพนธ์ประจำภาคการศึกษาที่ 2/2565
  \item นำเสนอปริญญานิพนธ์ประจำภาคการศึกษาที่ 2/2565
\end{enumerate}

\section{ประโยชน์ที่คาดว่าจะได้รับ}

โครงงานนี้จะเป็นประโยชน์กับใคร ยังไง ทั้งในเชิงรูปธรรมและนามธรรม ในปัจจุบันหรือในอนาคตถ้านำไป
ต่อยอด

\section{ตารางการดำเนินงาน}

\newpage
\section{ผลการดำเนินงาน}
ผลการดำเนินงานในภาคการศึกษาที่ 1/2565
\begin{enumerate}
  \item โครงสร้างและรูปแบบของเว็บแอปพลิเคชัน
  \item เก็บโครงสร้างของข้อมูลที่ได้รับจากอุปกรณ์เครือข่าย
  \item รูปเล่มปริญญานิพนธ์แสดงความคืบหน้าในภาคการศึกษาที่ 1/2565
\end{enumerate}
ผลการดำเนินงานในภาคการศึกษาที่ 2/2565
\begin{enumerate}
  \item เว็บแอปพลิเคชันที่ทำงานได้อย่างสมบูรณ์
  \item รูปเล่มปริญญานิพนธ์
\end{enumerate}

%%%%%%%%%%%%%%%%%%%%%%%%%%%%%%%%%%%%%%%%%%%%%%%%%%%%%%%%%%%%
%%%%%%%%%%%%%%%%%%%%%%% CHAPTER2 %%%%%%%%%%%%%%%%%%%%%%%%%%%
%%%%%%%%%%%%%%%%%%%%%%%%%%%%%%%%%%%%%%%%%%%%%%%%%%%%%%%%%%%%
\chapter{ทฤษฎีความรู้และงานที่เกี่ยวข้อง}

\hspace*{1cm}ในการจัดทําปริญญานิพนธ์เรื่อง Document management system for lawyer มีการอ้างอิงถึงทฤษฎี งานวิจัย และผลิตภัณฑ์ที่เกี่ยวข้องที่ใช้ในการพัฒนาชิ้นงาน ดังต่อไปนี้

\section{ทฤษฎีที่เกี่ยวข้อง}


\subsection{การดำเนินคดีของทนายความ} 
\hspace*{1cm}หลังจากที่ ทนายได้รับเรื่องคดีเข้ามาแล้ว ทนายก็จะทำการสอบข้อเท็จจริงจากลูกความ และทำการปรับข้อเท็จจริงให้เข้ากับหลักกฎหมาย จากนั้น จะทำการออกหนังสือกล่าวทวงถาม เขียนคำฟ้อง คำให้การ และคำให้การพร้อมฟ้องแย้ง หลังจากนั้นจะทำการยื่นคำฟ้อง หรือยื่นคำให้การให้แก่ศาล ต่อมาจะเป็นการ ยื่นบัญชีระบุพยาน และทำการขอหมายเรียก พยานเอกสาร พยานบุคคล หรือพยานวัตถุ จากนั้นจะเป็นการเตรียมคดี หรือซ้อมพยาน ก่อนที่จะนัดฟังคำพิพากษา ทั้งศาลชั้นต้น ศาลอุทธรณ์ และศาลฎีกา 

\subsection{Kanban software developments} 
\hspace*{1cm}Kanban board ถูกสร้างจากวิศวกรของบริษัท Toyota ที่มีชื่อว่า Taiichi Ohno ในปี 1940 โดนแนวคิดนี้สร้างขึ้นเพื่อควบคุมและจัดการงาน สินค้าคงคลังให้ในแต่ละขึ้นตอน ผลิตอย่างเหมาะสม และเป็นแบบแผนที่สุด\\
\hspace*{1cm}ดังนั้น Kanban board ถูกออกแบบมาเพื่อให้เห็นภาพรวม และช่วยควบคุม workflow ทั้งหมด ของโปรเจคนั้น ซึ่งจะมีลักษณะเป็น การดานหรือบอร์ด และการ์ด โดยในบอร์ดนั้น จะแสดงสเตตัสงานของทีม โดยจะมี post-it หรือการ์ด ที่จะมี Task หรืองานที่เราจะต้องทำ แปะอยู่ในสเตตัสนั้นๆ เพื่อทำให้เราสามารถรู้ได้ว่า เราควรจะต้องโฟกัสกำงานใดก่อน โดยสถานะพื้นฐานที่คนส่วนใหญ่ใช้กันก็จะมี To Do, Doing และDone แต่เราก็สามารถเพิ่มลดได้เองตามความเหมาะสมกับทีมและ เนื้องาน

\subsection{Database System}
\hspace*{1cm} Database หรือ ฐานข้อมูล คือ หมายถึง กลุ่มของข้อมูลที่ถูกเก็บรวบรวมไว้ โดยมีความสัมพันธ์กัน และถูกจัดเก็บอย่างมีระบบ เพื่อตอบสนองต่อความต้องการขององค์กร หรือระบบ โดยไม่จำเป็นต้องอยู่ในแฟ้มข้อมูลเดียวกัน สามารถแยกเก็บหลาย ๆ แฟ้มข้อมูลได้\\
\hspace*{1cm} Database System หรือ ระบบฐานข้อมูล หมายถึง การรวมกันของฐานข้อมูลที่เกี่ยวข้องกันตั้งแต่ 2 ฐานข้อมูลขึ้นไป โดยมีจุดประสงค์เพื่อลดความซ้ำซ้อนของข้อมูล และช่วยเพิ่มประสิทธิภาพในการจัดการข้อมูล ผ่านระบบจัดการฐานข้อมูล (Database Management System) ซึ่งการแบ่งชนิดของฐานข้อมูลสามารถแบ่งได้หลายหัวข้อ แต่ในที่นี้ผู้จัดทำจะแบ่งชนิดของฐานข้อมูลตามความนิยมใช้งาน โดยจะสามารถแบ่งออกได้เป็น 3 ประเภท ดังนี้
\begin{enumerate}
    \item {RDBMS (Relational Database Management System)}\\
    \hspace*{1cm}เป็นระบบจัดการฐานข้อมูลที่มีความเสถียรมากเหมาะสำหรับการเก็บข้อมูลที่มีจุดประสงค์และแยกประเภทชัดเจน ใช้ภาษา SQL ในการ Query และ Maintan Database มีการเก็บข้อมูลอยู่ในรูปแบบ ตาราง (Table) มีองค์ประกอบเป็น Rows และ Columns ซึ่งอำนวยความสะดวกในการเข้าถึงข้อมูลที่มีความสัมพันธ์กันระหว่างตาราง
    \item {NoSQL Database}\\
    \hspace*{1cm}NoSQL หรือ Non-relational database เป็นระบบจัดการฐานข้อมูลที่ไม่มีความสัมพันธ์กันอย่างชัดเจนหรือไม่มีรูปแบบที่แน่นอน เหมาะสำหรับการใช้งานกับ Big Data Real-time Web Application หรือ Logging โดยให้ความสำคัญกับเรื่องของความเร็วมากกว่าความถูกต้องแม่นยำ รวมถึงเหมาะกับข้อมูลชนิด Dynamic Data ที่มีการเปลี่ยนแปลงเกิดขึ้นบ่อยโดยไม่สนใจความสัมพันธ์ของข้อมูลก่อนหน้า
    \item {In-memory Database}\\
    \hspace*{1cm}เป็นระบบจัดการฐานข้อมูลชนิดหนึ่งที่ถูกสร้างขึ้นเพื่อวัตถุประสงค์ในการเก็บข้อมูลที่ต้องการเข้าถึงอย่างรวดเร็ว โดยใช้วิธีการเก็บ data ไว้ใน memory แทน ซึ่งมีความเร็วมากกว่า การจัดเก็บข้อมูลประเภทอื่น เช่น HDD อย่างไรก็ตาม In-memory database นั้นมีความเสี่ยงที่จะเกิดการสูญหายของข้อมูลได้มาก ดังนั้น In-memory database จึงเหมาะสำหรับ Application ที่ต้องการความรวดเร็วในการตอบสนองอย่างมาก
\end{enumerate}

\subsection{RESTful API}
\hspace*{1cm}RESTful API เป็นรูปแบบของสถาปัตยกรรมที่ใช้ Web Protocol ในการติดต่อสื่อสาร ระหว่างคอมพิวเตอร์สองเครื่อง ซึ่งในการติดต่อสื่อสาร แต่ละครั้งต้องมีการกำหนดรูปแบบการใช้งานก่อน เช่น GET POST เป็นต้น ซึ่งความสามารถ หลักๆของ RESTful นั้นคือสามารถส่งข้อมูล เก็บข้อมูล แสดงผล โดยจะมีทั้ง รูปภาพ, วิดีโอ, หรือเป็นข้อมูลที่เกี่ยวกับธุรกิจ\\
\hspace*{1cm}ซึ่งใน RESTful API จะเป็นการส่ง HTTP Request ไปยังเครื่องเซิฟเวอร์ โดยใน HTTP Request นั้นจะประกอบไปด้วย\\
\begin{itemize}
  \item {VERB คือเป็น HTTP Method เช่น GET POST PATCH DELETE เป็นต้น}
  \item {URI คือ จุดหมายหรือตำแหน่งที่ต้องการให้ระบบทำงาน}
  \item {HTTP Version คือเวอร์ชั่นของ HTTP ว่า HTTP Request นี้ใช้เวอร์ชั่นอะไร}
  \item {Request Header คือ metadata ที่เอาไว้ใช้เก็บ key-header เพื่อระบุข้อมูลของผู้ส่ง ผ่านทาง Header}
  \item {Request Body คือข้อมูล ในส่วนของ RESTful ที่เราจะทำการส่ง หรือได้รับ ซึ่งสามารถประกอบไปด้วย รูปภาพ วิดีโอ หรือข้อมูลที่เกี่ยวกับธุรกิจ}  
\end{itemize}

\subsection{Data visualization}

\subsection{Optical Character Recognition}

\newpage
\section{งานวิจัยที่เกี่ยวข้อง}

\section{ผลิตภัณฑ์ที่ที่เกี่ยวข้อง}
\subsection{Avokaado application}
\hspace*{1cm}Avokaado คือ all-in-one contract lifecycle management ที่จะช่วยสำนักงานกฎหมาย ในการสร้าง จัดการ และดูแลเอกสารต่างๆ อยู่ในแพลตฟอร์มเดียว เช่น smart drafting, co-editing, negotiating approving, e-signing, document management และ tracking dashboard
% \begin{figure}[!h]\centering
    %   \includegraphics[width=13cm]{./image/NetScout.png}
    %   \caption{ รูปแสดงซอฟแวร์ nGeniusPLUSE ของ NETSCOUT }\label{fig:netScout}
    % \end{figure}
\subsection{ContractSafe application}
\hspace*{1cm}xxx

\section{เครื่องมือที่ใช้ในการพัฒนา}

%%%%%%%%%%%%%%%%%%%%%%%%%%%%%%%%%%%%%%%%%%%%%%%%%%%%%%%%%%%%
%%%%%%%%%%%%%%%%%%%%%%% CHAPTER3 %%%%%%%%%%%%%%%%%%%%%%%%%%%
%%%%%%%%%%%%%%%%%%%%%%%%%%%%%%%%%%%%%%%%%%%%%%%%%%%%%%%%%%%%
\chapter{วิธีการดำเนินงานกระบวนการและการออกแบบ}

Explain the design (how you plan to implement your work) of your project. Adjust the section titles below to suit the types of your work. Detailed physical design like circuits and source codes should be placed in the appendix.

\section{ข้อกำหนดและความต้องการของระบบ}


\section{Hardware Module 1}
\subsection{Component 1}
\subsection{Logical Circuit Diagram}

\section{Hardware Module 2}
\subsection{Component 1}
\subsection{Component 2}

\section{Path Finding Algorithm}

\section{Database Design}

\section{UML Design}

\section{GUI Design}

\section{การออกแบบการทดลอง}
\subsection{ตัวชี้วัดและปัจจัยที่ศึกษา}
\subsection{รูปแบบการเก็บข้อมูล}




%%%%%%%%%%%%%%%%%%%%%%%%%%%%%%%%%%%%%%%%%%%%%%%%%%%%%%%%%%%%
%%%%%%%%%%%%%%%%%%%%%%% CHAPTER4 %%%%%%%%%%%%%%%%%%%%%%%%%%%
%%%%%%%%%%%%%%%%%%%%%%%%%%%%%%%%%%%%%%%%%%%%%%%%%%%%%%%%%%%%
\chapter{ผลการดำเนินงาน}

You can title this chapter as \textbf{Preliminary Results} ผลการดำเนินงานเบื้องต้น or \textbf{Work Progress} ความก้าวหน้าโครงงาน for the progress reports. Present implementation or experimental results here and discuss them.
ใส่เฉพาะหัวข้อที่เกี่ยวข้องกับงานที่ทำ 

\section{ประสิทฺธิภาพการทำงานของระบบ} 
\section{ความพึงพอใจการใช้งาน}
\section{การวิเคราะห์ข้อมูลและผลการทดลอง}

%%%%%%%%%%%%%%%%%%%%%%%%%%%%%%%%%%%%%%%%%%%%%%%%%%%%%%%%%%%%
%%%%%%%%%%%%%%%%%%%%%%% CHAPTER5 %%%%%%%%%%%%%%%%%%%%%%%%%%%
%%%%%%%%%%%%%%%%%%%%%%%%%%%%%%%%%%%%%%%%%%%%%%%%%%%%%%%%%%%%
\chapter{บทสรุป}

This chapter is optional for proposal and progress reports but 
is required for the final report.

\section{สรุปผลโครงงาน}
สรุปว่าโครงงานบรรลุตามวัตถุประสงค์ที่ตั้งไว้หรือไม่ อย่างไร 

\section{ปัญหาที่พบและการแก้ไข}
State your problems and how you fixed them.

\section{ข้อจำกัดและข้อเสนอแนะ}
ข้อจำกัดของโครงงาน What could be done in the future to make your projects better.

%%%%%%%%%%%%%%%%%%%%%%%%%%%%%%%%%%%%%%%%%%%%%%%%%%%%%%%%%%%%%%%
%%%%%%%%%%%%%%%%%%%% Bibliography %%%%%%%%%%%%%%%%%%%%%%%%%%%%%
%%%%%%%%%%%%%%%%%%%%%%%%%%%%%%%%%%%%%%%%%%%%%%%%%%%%%%%%%%%%%%%

%%%% Comment this in your report to show only references you have
%%%% cited. Otherwise, all the references below will be shown.
%\nocite{*}
%% Use the kmutt.bst for bibtex bibliography style 
%% You must have cpe.bib and string.bib in your current directory.
%% You may go to file .bbl to manually edit the bib items.

\makeatletter
\g@addto@macro{\UrlBreaks}{\UrlOrds}
\makeatother

\bibliographystyle{kmutt}
\bibliography{string,cpe}

%%%%%%%%%%%%%%%%%%%%%%%%%%%%%%%%%%%%%%%%%%%%%%%%%%%%%%%%%%%%%%%
%%%%%%%%%%%%%%%%%%%%%%%% Appendix %%%%%%%%%%%%%%%%%%%%%%%%%%%%%
%%%%%%%%%%%%%%%%%%%%%%%%%%%%%%%%%%%%%%%%%%%%%%%%%%%%%%%%%%%%%%%
\appendix{ชื่อภาคผนวกที่ 1}
\noindent{\large\bf ใส่หัวข้อตามความเหมาะสม} \\





\end{document}
